\section{The verification script}\label{app:scripts}

The output below is the first part of a run of \texttt{verify\_all.py}, and
covers items (I) to (V) of \Cref{sec:verification}; the remaining items are
carried by the companion scripts named there, and the whole run ends with the
line \texttt{947 checks passed, 0 failed}. Everything requires the standard library of Python~3,
together with \texttt{sympy} for \texttt{secant\_plane.py} and
\texttt{weiltype\_family.py}, where it does exact symbolic linear algebra over
$\QQ(\sqrt{-d})$, for \texttt{mumford\_rigidity.py}, where it factors
Gram determinants over $\QQ[a_{0},\dots,a_{3}]$, and for
\texttt{mumford\_object.py}, \texttt{lefschetz\_closure.py},
\texttt{twistor\_locus.py} and \texttt{hk\_pullback.py}, where it computes
exact ranks, null spaces, determinants and characteristic polynomials, and
\texttt{numpy} for \texttt{pte\_remaining.py} and \texttt{weil\_tori.py},
where it holds integer arrays and reduces them modulo a prime; no step uses
floating point.

\begin{small}
\begin{verbatim}
(I) the Fourier-Mukai transform on powers of the polarisation
  [PASS] g=1: every image is a pure multiple of theta^(g-k)
  [PASS] g=2: every image is a pure multiple of theta^(g-k)
         k=0: -1/2  k=1: 1  k=2: -2
  [PASS] g=3: multiples agree with the closed form
         k=0: 1/6  k=1: -1/2  k=2: 2  k=3: -6
  [PASS] g=4: multiples agree with the closed form
         k=0: 1/24  k=1: -1/6  k=2: 1  k=3: -6  k=4: 24

(II) the sign sigma(g,k), uniform over index sets
  [PASS] sign is the same for every index set T, all g <= 12
  [PASS] sign equals (-1)^(g(g+1)/2 + k), all g <= 12

(III) the character grading: Weil line against eta^n
  [PASS] Weil line and eta^n occupy different summands, 1 <= n <= 4
         n=2: Weil line at (r,s) = [(4,0), (0,4)],  eta^2 at (2,2)
         n=3: Weil line at (r,s) = [(6,0), (0,6)],  eta^3 at (3,3)
  [PASS] the characters tau^(2n) and N^n are never equal for n >= 1

(IV) the secant count against the Lefschetz thresholds
  [PASS] 2n^2+2n < 2^(2n) < 3^(2n) for 2 <= n <= 8
         n=2: chi required 12,  2^(2n) 16,  3^(2n) 81
         n=3: chi required 24,  2^(2n) 64,  3^(2n) 729
  [PASS] counterexample: (1,3) is base point free with chi = 3 < 4
  [PASS] counterexample: (1,5) is very ample with chi = 5 < 9
  [PASS] so neither threshold is a necessary condition

(V) the semiregularity target
  [PASS] sum_q h^(q,q+2) = C(4n,2n-2) for 2 <= n <= 7
         n=2: 28    n=3: 495   n=4: 8008   n=5: 125970
  [PASS] the target grows by a factor above 10 at each step

  18 checks passed, 0 failed
  overall: PASS
\end{verbatim}
\end{small}

Item (XXI), the existence of a sheaf with the required Chern character, is the
one item whose output is an explicit witness rather than a verdict, so we
reproduce the dimension four block of \texttt{secant\_exists.py} as well.

\begin{small}
\begin{verbatim}
  == n = 4 ==
       d=3, a=1, b=1:  M = 6,
       alpha = -23[O(0t)] + 66[O(1t)] - 54[O(2t)] + 20[O(3t)] - 3[O(4t)]
       rank alpha = 6 = M a
  [PASS] n=4: the Vandermonde system is solvable for every (d,a,b)
  [PASS] n=4: every witness reproduces M c_k in every degree
  [PASS] n=4: every witness has rank exactly M a > 0
  [PASS] n=4: the multiple M is bounded over the whole range
         largest M = 12, and M = 1 in 4 of 72 cases
  [PASS] n=4: theta^k/k! has integer coefficients in a symplectic basis
\end{verbatim}
\end{small}

Item (XXXIII), the closure graph, is the other item whose output is worth
reading rather than merely passing, since it is the statement of what is left.
The block below is the frontier it computes, over the rules of this paper and
those of \Cref{ssec:bullseye}, followed by the derivation it prints for each
minimal sufficient set.

\begin{small}
\begin{verbatim}
  [PASS] the minimal sufficient sets are enumerated exhaustively
         5 of them, over the 16 open statements, from 65536 subsets
  [PASS] the minimal sufficient sets are exactly five
         {red_ab}, {lef_B, mot}, {lef_B, red_ab_mod}, {red_ab_mod, vhc}
         and {P2_cm_s, red_ab_mod, red_weil}
  [PASS] exactly one single open statement suffices, and it is red_ab
         the conjecture for the varieties that are not abelian gives the
         conjecture for abelian varieties through A x P^1, and then the
         conjecture
  [PASS] the route through this paper and the two routes of the
         literature with red_ab suffice and are not minimal
         each contains the sufficient singleton {red_ab}; with
         red_ab_mod in its place each is minimal
  [PASS] red_ab_mod alone gives neither the conjecture for abelian
         varieties nor the conjecture
         it holds trivially on abelian varieties, so it needs a second
         statement that reaches them
  [PASS] the route through this paper, with red_ab_mod, is the only
         minimal sufficient set with three elements
         it needs one statement more than each route of the literature
         through red_ab_mod

       the minimal sufficient sets:
            {red_ab}
            {lef_B, mot}
            {lef_B, red_ab_mod}
            {red_ab_mod, vhc}
            {P2_cm_s, red_ab_mod, red_weil}
       with
            P2_cm_s        (P2) for the split Weil families of the CM
                           fields of degree at least four
            lef_B          the Lefschetz standard conjecture B(X) for
                           every smooth complex projective variety X
            mot            every Hodge class on every smooth complex
                           projective variety is motivated in the sense
                           of Andre
            red_ab         the Hodge conjecture for every smooth
                           projective variety that is not an abelian
                           variety; equivalent to the conjecture, since
                           it covers A x P^1 for every abelian variety A
                           (prop:f3ishc)
            red_ab_mod     every Hodge class on every smooth projective
                           variety is algebraic modulo images of Hodge
                           classes of abelian varieties under algebraic
                           correspondences (def:f3prime); trivial on
                           abelian varieties
            red_weil       every Hodge class on an abelian variety
                           outside the subring generated by divisor
                           classes and Weil classes of quotients is
                           algebraic; not a reduction, since the self-
                           product of a Mumford fourfold carries two
                           such classes
            vhc            a global section of R^{2p} f_* Q for a smooth
                           projective family over a smooth connected
                           base that is algebraic at one point is
                           algebraic at every point

  [PASS] every member of a minimal sufficient set other than P2_cm_s is
         a consequence of the conjecture
         lef_B by thm:equivalence, mot by prop:lefmot, vhc by
         prop:lefvhc, red_ab_mod by prop:f3prime, red_ab and red_weil
         directly; no implication from the conjecture to P2 is known
  [PASS] {red_ab} and {lef_B, mot} are equivalences with the conjecture
         {red_ab} by prop:f3ishc, {lef_B, mot} by prop:lefmot
  [PASS] every minimal sufficient set contains red_ab, red_ab_mod or mot
         the varieties that are not abelian are reached only through the
         conjecture for them, through the conjecture modulo abelian
         varieties, or through motivated classes
  [PASS] the variational statement alone gives the conjecture for
         abelian varieties and not the conjecture
         so it covers the Weil classes of every CM field and the classes
         beyond the Weil lines at once
  [PASS] the Lefschetz standard conjecture gives the variational
         statement
         through the deformation theorem for motivated classes
  [PASS] only the route through this paper uses a statement proved in it
         the four other minimal sets rest on open statements and quoted
         theorems alone
  [PASS] each element of a minimal sufficient set is necessary to it
         removing any one of them leaves the conjecture underivable
  [PASS] the secant route lies in no minimal sufficient set
         granting all six of its open demands does not put it on a path
         to the conjecture
  [PASS] the secant route, fully granted, reaches every imaginary
         quadratic family and stops short of the fields of higher degree
         its base points are split members, of trivial discriminant, and
         descent (prop:descent) carries the trivial discriminant to
         every other; it has no base point over a field of degree at
         least four
  [PASS] the trivial discriminant closes the imaginary quadratic case
         by descent: B x Y, with Y a Weil surface of any discriminant,
         and a push-forward against the Weil class of Y
  [PASS] the split families of the fields of degree at least four close
         the imaginary quadratic case as well
         by scalar extension B -> B (x) O_F, every K lying in K(sqrt 2)
  [PASS] the propagation statement for the split imaginary quadratic
         families closes the imaginary quadratic case outright
         a base point exists in every split family, and descent does the
         rest
  [PASS] the propagation statement for the split families of the fields
         of degree at least four closes the Weil classes of every CM
         field
         descent within each field, and scalar extension down to the
         imaginary quadratic fields
  [PASS] the families of nontrivial discriminant need no propagation
         statement of their own
         no minimal sufficient set contains (P2) for a family of
         nontrivial discriminant or for an imaginary quadratic field
  [PASS] the bounded criterion (P1) is not in the rule set
         it is equivalent to the conjecture for the family, so it is a
         restatement and not a premise
  [PASS] the Weil classes of the higher CM fields are derived, not open
  [PASS] without prop:f3ishc and red_ab_mod the search returns the four
         sets of the earlier version
         {lef_B, mot}, {lef_B, red_ab}, {red_ab, vhc}, {P2_cm_s, red_ab,
         red_weil}: the change is those two additions, with (P2) read on
         the split families of the higher fields
  [PASS] the rule of prop:f3ishc has red_ab as its only premise
         and the conjecture for abelian varieties as its conclusion: the
         conjecture for A x P^1 gives it for A

       the conjecture from the proved and quoted statements and red_ab
            HC_ab         from red_ab
                                                             [prop:f3ishc]
            HC            from HC_ab, red_ab
                                                         [ssec:closuremap]

       the conjecture from the proved and quoted statements and lef_B,
         mot
            HC            from lef_B, mot
                                                             [prop:lefmot]

       the conjecture from the proved and quoted statements and lef_B,
         red_ab_mod
            vhc           from lef_B, mot_def
                                                             [prop:lefvhc]
            HC_ab         from vhc, acc_ab
                                                               [thm:vhcab]
            HC            from HC_ab, red_ab_mod
                                                            [prop:f3prime]

       the conjecture from the proved and quoted statements and
         red_ab_mod, vhc
            HC_ab         from vhc, acc_ab
                                                               [thm:vhcab]
            HC            from HC_ab, red_ab_mod
                                                            [prop:f3prime]

       the conjecture from the proved and quoted statements and P2_cm_s,
         red_ab_mod, red_weil
            weil_cm_triv  from base_point_cm, reduction_cm,
                               orbit_dense_cm, P2_cm_s
                                                       [thm:cmpropagation]
            weil_cm_nontriv from weil_cm_triv
                                                            [prop:descent]
            weil_triv     from weil_cm_triv
                                                            [prop:descent]
            weil_cm       from weil_cm_triv, weil_cm_nontriv
                                                             [prop:cmweil]
            weil_nontriv  from weil_triv
                                                            [prop:descent]
            weil_iq       from weil_triv, weil_nontriv
                                                           [prop:separate]
            weil_all      from weil_iq, weil_cm
                                                              [app:albert]
            HC_ab         from weil_all, red_weil
                                                         [ssec:closuremap]
            HC            from HC_ab, red_ab_mod
                                                            [prop:f3prime]
\end{verbatim}
\end{small}

Item (XLI) prints the numbers on which \Cref{thm:mumfordks} turns: the
scalar $\nu$, the transport of the real multiplication, and the shape of the
Kuga--Satake map on the explicit Clifford algebra.

\begin{small}
\begin{verbatim}
  [PASS] mu mu^dagger acts on each U_ij by one and the same nonzero
         rational nu
         nu = -1 on U_12, U_13, U_23; 27/2 on the line of theta
  [PASS] with kappa = e and the form q_lambda, mu mu^dagger acts on U_ij
         by nu g_i g_j, where g = e^2/lambda
         e = (2, -1, 5), lambda = (3, 7, 2): on U_ij by nu g_i g_j,
         g = (4/3, 1/7, 25/2)
  [PASS] kappa^dagger kappa, for the trace form on End(V), is
         multiplication by 4 e^2/lambda
  [PASS] C^+(T) has dimension 256 and is V^{32}: the vectors of highest
         weight (1,1,1) span 32 dimensions, and lowering them gives a basis
         dim C^+ = 256, highest weight vectors 32
  [PASS] N_1, N_2, N_3 anticommute pairwise and each squares to a nonzero
         scalar, proportional to the conjugate of lambda
         N_i^2 = 4, 8, 20 for lambda = (1, 2, 5)
\end{verbatim}
\end{small}

Item (XLIV) prints the checks behind \Cref{thm:mumfordrigid} and
\Cref{thm:mumfordnumerical}: the weight blocks, the factors of their Gram
determinants, the one kernel vector, and the certificates.

\begin{small}
\begin{verbatim}
  [PASS] wedge^2 V splits as 1 + 9 + 9 + 9 and the four tensors pi are
         invariant
         dimensions {'0': 1, '12': 9, '13': 9, '23': 9}
  [PASS] sp^{-1,1} has dimension 36, split into nine weight blocks
         block dimensions {(-2, -2): 3, (-2, 0): 4, (-2, 2): 3, (0,
         -2): 4, (0, 0): 8, (0, 2): 4, (2, -2): 3, (2, 0): 4, (2, 2):
         3}
  [PASS] the generic kernel is one vector, in the block of weight (0,0),
         and it does not depend on a
         the tangent direction of the Mumford curve
  [PASS] every Gram determinant is a positive constant times a product of
         P1, P2, Q1, Q2, Q3, Q4 and a2^2 + a3^2
  [PASS] each of the seven forms is a sum of squares of rational linear
         forms with positive coefficients
  [PASS] the common zeros of the squared linear forms of each form have
         two of a1, a2, a3 equal
         so for pairwise distinct a1, a2, a3 every block has full rank
         and the tangent space is one dimensional
  [PASS] tangent dimensions: 1 for distinct coefficients, 10 for a1 = a2
         = a3, 20 for a0 = a1 = a2 = a3
         computed 1, 10, 20
  [PASS] the annihilator in HH^2 (dimension 120) of an exceptional class
         is one dimensional at forty random admissible points, and 36
         dimensional for the class with equal coefficients
         so on a nonempty Zariski open set of classes, which contains
         rational ones, the least dimension of Ext^2 is 119
  [PASS] the annihilator does jump on special loci, which contain no
         rational exceptional class
         dimension 7 at a2 = 0 and 4 on the hyperplane a0 - 3 a1 - 3 a2
         + a3 = 0
  [PASS] HH^2 = H^{0,2} + H^1(T) + H^0(wedge^2 T) splits into blocks by
         type, weight and exchange parity, and omega_a is exchange
         invariant
         dimensions {'pi': 28, 'v': 64, 'z': 28}
  [PASS] every block Gram determinant is a positive constant times a
         product of P1, P2, Q1..Q4, R, S1..S4, Z1..Z6, a2, a3, L0, L1,
         L2, L3
  [PASS] the only block with a kernel at a generic point is the exchange
         even block of weight (0,0) of H^1(T), its kernel is independent
         of a and is the derivation of the Mumford curve
  [PASS] L0 = a0 + a1 + a2 + a3 divides only the Gram determinants of the
         exchange odd blocks of bivectors
         9 blocks
  [PASS] S1..S4 are sums of squares whose zeros force two of a1, a2, a3
         to coincide or one to vanish; Z1..Z6 are positive definite, as
         sums of four squares of independent forms with positive
         coefficients
  [PASS] every linear factor other than L0 has unequal coefficients on
         a1, a2, a3
         so it does not vanish at the conjugates of an irrational
         element of a cubic field, whatever the rational a0
  [PASS] on the hyperplane L0 = 0 the annihilator has dimension 10
         at three sample points: [10, 10, 10]
\end{verbatim}
\end{small}

Item (XLV) prints the checks behind \Cref{thm:mumfordprofile} and
\Cref{prop:mumforddivisor}.

\begin{small}
\begin{verbatim}
  [PASS] the ranks of contraction into omega at (3,5,-7,11), exact over
         Q, are 1,16,119,328,560,328,119,16,1 in degrees 0..8
         computed [1, 16, 119, 328, 560, 328, 119, 16, 1]
  [PASS] they vanish in degrees above eight
  [PASS] the same ranks, modulo p = 1000133, at a0 = 2 and the conjugates
         of 2cos(2pi/7) in each of the six orders
         a rational exceptional class for the cubic subfield of
         Q(zeta_7)
  [PASS] on the hyperplane a0 + a1 + a2 + a3 = 0 the ranks drop to
         1,16,110,328,548,328,110,16,1
         computed [1, 16, 110, 328, 548, 328, 110, 16, 1]
  [PASS] for a random class with components of types (0,0) to (3,3) the
         ranks are palindromic, r_k = r_{8-k}, and zero above eight
         computed [1, 16, 104, 322, 474, 322, 104, 16, 1, 0, 0]
  [PASS] the alternating sum of the ranks is 112, and 2r_0 + 2r_2 + r_4 =
         800
         so chi(E,E) = 0 forces Ext^1 + Ext^3 >= 400 when Ext^{<0} = 0
  [PASS] at a CM point: 16 divisor classes, 132 Hodge classes in degree
         four, products of divisor classes spanning 100
  [PASS] the combinations of the pi in that span are exactly those with
         a1 = a2 = a3
         so no rational exceptional class is a polynomial in divisor
         classes at any point
  [PASS] at a CM point X_c has 4 divisor classes and 8 Hodge classes in
         degree four, of which products of divisors span 6
         the other two are not products of divisor classes
\end{verbatim}
\end{small}

Item (XLVI) prints the checks behind \Cref{thm:lefschetzclosure}.

\begin{small}
\begin{verbatim}
  [PASS] the 36 maps lie in sp(V, psi), there are 32 root vectors, and
         the Lie algebra of G lies in sp(V, psi)
  [PASS] the Sp(V, psi)-invariants of wedge^k(V+V) have dimension
         1,3,6,10,15,10,6,3,1, and the products of the three invariant
         divisor classes span them
         upper [1, 3, 6, 10, 15, 10, 6, 3, 1],
         lower [1, 3, 6, 10, 15, 10, 6, 3, 1]
  [PASS] the G-invariants have dimension 1,3,8,16,28,16,8,3,1 (by
         weight multiplicities, and as the kernel modulo p), so the
         Hodge classes that are not Lefschetz number 0,0,2,6,13,6,2,0,0,
         29 in all
         computed [1, 3, 8, 16, 28, 16, 8, 3, 1]
  [PASS] at a CM point the T_Sp-invariant monomials number
         1,16,100,304,454,304,100,16,1 = coefficients of (1+4y+y^2)^4,
         each a product of the 16 divisor classes, which span exactly
         them
         Hodge classes [1, 16, 132, 432, 648, 432, 132, 16, 1]
  [PASS] at a CM point the Hodge classes number
         1,16,132,432,648,432,132,16,1
  [PASS] pi_0 and pi_12 + pi_13 + pi_23 are Sp-invariant and pi_12 is
         not; omega_a is T_Sp-invariant exactly when a1 = a2 = a3
         12 monomials of nonzero T_Sp-weight, exact rank 2
  [PASS] the module generated by pi_12 - pi_13, and the one generated by
         pi_13 - pi_23, has dimension 308 and highest weight (2,2,0,0)
         dimensions [308, 308]
  [PASS] Weyl's formula for Sp_8: dimensions 1, 27, 42, 308 for 0, w_2,
         w_4, 2w_2, and 1 + 27 + 42 + 308 = 378 = 27*28/2
  [PASS] that module meets the span of the four pi in a plane: the
         exceptional plane
         computed [2, 2]
\end{verbatim}
\end{small}

Item (XLVII) prints the checks behind \Cref{prop:notwistor}.

\begin{small}
\begin{verbatim}
  [PASS] s = s_1 (x) s_2 (x) s_3 is a real structure, V_R has real
         dimension 8, and psi is real and alternating on V_R
  [PASS] J_3, J_1, J_2 preserve V_R and square to -1; psi(x, Jy) is real
         symmetric of signature (0,8) for J_3 and (4,4) for J_1 and J_2
         signatures {'J_3': (0, 8, 0), 'J_1': (4, 4, 0),
         'J_2': (4, 4, 0)}
  [PASS] at J_1 the quaternion k = s_1 (x) 1 (x) 1 preserves V_R and
         psi, anticommutes with J_1 and carries psi(x, J_1 y) to its
         negative; so psi (x) S is congruent to its negative, hence
         indefinite, for every nonzero real symmetric 2 x 2 matrix S: no
         invariant class of degree two polarises the square
         signature (m,m) confirmed directly on 6 matrices S
  [PASS] at J_1 the tensors pi have type (2,2) and the annihilator of
         omega_a in H^1(T) (dimension 64) is spanned by the lowering
         derivation of the first factor on both copies
         dimension 64, kernel 1
  [PASS] at J_3, the Mumford family, it is spanned by the lowering
         derivation of the third factor, as in item (XLIV)
         dimension 64, kernel 1
  [PASS] the exchange of the first and third factors carries pi_12 to
         pi_23 and fixes pi_0 and pi_13
\end{verbatim}
\end{small}

Item (XLVIII) prints the checks behind \Cref{prop:hksquare}.

\begin{small}
\begin{verbatim}
  [PASS] iota(x) = (x (x) 1) Psi is symmetric for the nine generators of
         T, and y.iota(x) = iota([y, x]) for all pairs of generators
  [PASS] the Casimirs split H^2 of the square (dimension 120) into
         isotypic pieces of dimensions 3, 27, 27, 27, 27, 3, 3, 3 with
         (2,0)-parts 0, 0, 9, 9, 9, 0, 0, 1; the pieces of types
         (2,0,0), (0,2,0), (0,0,2) are iota(T_1), iota(T_2), iota(T_3),
         and only the Galois orbit of T has h^{2,0} = 1
         dims [3, 27, 27, 27, 27, 3, 3, 3]
  [PASS] sigma_0 = iota(E_3) is of type (2,0) and pairs the
         (1,0)-vectors of the two copies through an invertible 4 x 4
         block, so it is a holomorphic symplectic form on the square
         det of the block 1
  [PASS] iota_2(C_1) = 3 pi_0 - pi_12 - pi_13 + 3 pi_23 and cyclically;
         so iota_2(sum s_i C_i) = omega_a with a_0 = 3 sum s_i and
         a_ij = 4 s_k - sum s_i
  [PASS] det(x_1 + x_2 + x_3 on V) = Delta(N_1, N_2, N_3)^2 with N_i =
         -det(x_i), as a polynomial identity in nine variables
  [PASS] iota(x)^8 = +- 8! det(x) vol at sample points, with one sign
         ratios [1, 1, 1]
  [PASS] Delta is a nondegenerate ternary quadratic form (Gram
         determinant -4) and irreducible, so Delta(N)^2 is not a
         constant multiple of the fourth power of a linear form
         Gram determinant -4
\end{verbatim}
\end{small}

Item (XLIX) prints the checks behind \Cref{cor:mumfordone} and
\Cref{rem:mumfordneeded}.

\begin{small}
\begin{verbatim}
  [PASS] every theorem label in the rule set is a label of the paper,
         and every statement named in a rule is declared
         22 rules, 11 labels
  [PASS] the target is not in the closure of the proved statements
  [PASS] the statements equivalent to the target over the proved ones
         are exactly TARGET, UNCOUNT, B_WW, BOUND and MODP
         ['BOUND', 'B_WW', 'MODP', 'TARGET', 'UNCOUNT']
  [PASS] every minimal set of open statements yielding the target has
         one element, so one route suffices
         12 minimal sets: BOUND, B_WW, EXT119, F2, F3, HC, HK_U, KS_U,
         L, MODP, UNCOUNT, V
  [PASS] the target yields none of EXT119, KS_U, HK_U: those routes are
         stronger than the target, not forms of it
  [PASS] L and V each yield the target, so on the routes to the
         conjecture through them the target is not an input
  [PASS] HC and F2 yield the target, and the target yields none of HC,
         F1, F2, F3, L, V, M
  [PASS] F3 is equivalent to HC and is itself a one-element route to the
         target
         by prop:f3ishc; so the three statements of cor:fourremain are
         not three obligations either
  [PASS] the minimal sets of open statements yielding HC in this rule
         set are {HC}, {F3} and {L, M}, and none contains the target
         {F3}, {HC}, {L, M}
\end{verbatim}
\end{small}

The source is distributed with the paper as \texttt{verify\_all.py}. Its four
substantive routines are an exterior algebra over $\QQ$ with basis elements
indexed by bit masks and the shuffle sign \eqref{eq:shufflesign}; the
transform $\Phi_{\cP}$ implemented as multiplication by $e^{\ell}$ followed by
extraction of the coefficient of the top form in the dual variables; the
permutation sign of \Cref{lem:appsign}; and the binomial identities of
\Cref{prop:target}. No step uses floating point.
